\documentclass[12pt,addpoints]{exam}

\usepackage{mystyle}
\newcommand{\bCoord}[2]{\left[#1\right]_{#2}}
\newcommand{\Cb}{\mathcal{C}}
\renewcommand{\P}[1]{\mathcal{P}_{#1}\left(\mathbb{R}\right)}
\renewcommand{\B}{\mathcal{B}}
\usepackage{listings}

\title{MATH 3191-E01 Final Exam}
\date{Due 8 May 2026}

\author{Name: \rule{150pt}{1pt}}

\begin{document}
   \maketitle

    \paragraph{Instructions:} Submit ALL work that supports your answers. A solution with no work will
    receive a $0$. You have until 11:59pm (Mountain Time) Friday, May 8th to submit this assignment.
    I will close submissions on Canvas 2 hours after the due date to allow for tech issues. I will not
    accept late submissions without documentation supporting a university excused absence. If you are
    unsure, see \href{https://www.ucdenver.edu/docs/librariesprovider284/default-document-library/7030d---student-attendance-and-absences---cu-denver02fbb186-4b01-4e15-af21-2ab4bbaac065.pdf?sfvrsn=f1001eb4_1}
    {Policy File 7030D} for clarification. You are allowed to use any resource we have covered in
    class.
    I will only be using content from the chapters I have talked about in the slides. You can use content from the
    textbook that I haven't discussed (as long as it's in the above list!) and I won't count off,
    but there is another way to do it that may or may not be easier! Any homework assignment
    is allowed for reference.

    Also complete the file ``gram\_schmidt\_FIXME.ipynb''

    The written component (this document) has 170 points available (scored out of 160) while the
    jupyter notebook file is worth 20 points. So a total score of 190/180 is possible.


    \vfill
    \begin{center}
        \gradetable[h]
    \end{center}
    \newpage
    \begin{questions}
        \section*{``Proof'' Based Questions}
        \question[20] Let $A\in\R^{m\times n}$. In class, we discussed the SVD, which means we have an
        $U\in\R^{m\times m},\Sigma\in\R^{m\times n},$ and $V\in\R^{n\times n}$ such that
        \[
            A = U\Sigma V^\top
        \]

        Use this decomposition to prove that $\rowS{A}^\perp = \nullS{A}$.

        Hint: We gave a way to describe each of the four fundamental subspaces in the lecture notes.
        \vfill
        \newpage
        \question[20] Let $A\in\R^{m\times m}$ such that $\det{A} = 0$. Argue that $A$ must have
        linearly dependent rows/columns without using the invertible matrix theorem.
        Note: Consider $m=2,m=3$, then try to extend your argument.

        Hint: Consider how row operations affect the determinant
        \vfill
        \newpage
        \question[20] Let $A\in\R^{m\times n}$ and $\vec{b}\in\R^{m}$.
        Recall the least squares problem is to find an $\vec{x}\in\R^{n}$ such that
        \[
            \norm{A\vec{x} - b}{2}
        \]
        is minimal.

        Let $Q\in\R^{m\times m},\R^{m\times n}$ be the QR factorization of $A$. Demonstrate that
        the solution to
        \[
            A^\top A \vec{x} = A^\top\vec{b}
        \]
        and
        \[
            R \vec{x} = Q^\top\vec{b}
        \]
        both minimize
        \[
            \norm{A\vec{x} - b}{2}
        \]

        Hint: Show the norms of $ A^\top A \vec{x} = A^\top\vec{b} $ and $ R \vec{x} = Q^\top\vec{b} $
        are the same and use the fact that $\norm{Q}{2} = 1$
        \vfill
        \newpage
        \question[20] Let $V=\mathcal{P}_2(\R)$ be the set of all polynomials of degree at most $2$
        with real coefficients. Let $p = a_0 + a_1t + a_2t^2$ and $q = b_0 + b_1t + b_2t^2$ be elements
        of $V$

        \begin{enumerate}
            \item(10 points)
                Prove that
                \[
                    \langle p,q\rangle = a_0b_0 + a_1b_1 + a_2b_2
                \]
                is an inner product.
            \item(10 points)
                Compare this inner product to the dot product over $\R^3$
        \end{enumerate}
        \vfill
        \newpage
        \section*{Computational Questions}
        \question[20] Consider the following vector space
        \[
            V = \set{A\in\R^{2\times 2}}{\,\trace\left(A\right) = 0}
        \]
        Consider the basis
        \[
            \mathcal{B} = \setBasic{\bMat{1 & 0 \\ 0 & -1}, \bMat{0 & 1 \\ 0 & 0}, \bMat{0 & 0 \\ 1 & 0}}
        \]
        \begin{enumerate}
            \item(10 points) Give the $\mathcal{B}$ coordinates of each element in the set $S$ below
                \[
                    S = \setBasic{\bMat{2 & 1 \\ 0 & -2}, \bMat{1 & 1 \\ 0 & -1}, \bMat{-1 & -1 \\ 1 & 1}}
                \]
            \item(10 points) Using these coordinate vectors, determine if $S$ is a basis of $V$ as well
        \end{enumerate}
        \vfill
        \newpage
        \question[20] Let $V=\R^3$ be the vector space of arrays of real numbers of length 3. Consider
        the following basis of $V$
        \[
            \mathcal{C} = \setBasic{\bMat{1 \\ 1 \\ 2}, \bMat{0 \\ 2 \\ 1 }, \bMat{1 \\ 0 \\ 1}}
        \]
        Apply the Gram-Schmidt process to orthonormalize $\mathcal{C}$.
        \vfill
        \newpage
        \question[20] Let $A\in\R^{3\times 3}$ be given below
        \[
            A = \bMat{  4 & -2 & 0 \\
                        -2 & 0 & 1 \\
                        0 & 1 & 4}
        \]
        Diagonalize $A$. IE compute $V,D\in\R^{3\times 3}$ such that $D$ is diagonal and $A = VDV^{-1}$
        \vfill
        \newpage
        \question[20] Let $W = \Span{\bMat{1 \\1 \\ -1}, \bMat{0 \\ 1 \\ 1}}$. Find the orthogonal
        decomposition of $\vec{x} = \bMat{3 \\ 3 \\ -1}$.
        \vfill
        \newpage
        \section*{Bonus Question}
        \question[10] Answer one of the following questions for $10$ points. If you answer both you
        only receive credit for one! (Any answer gets full credit, have fun with it!)
        \begin{enumerate}
            \item Define the slang term ``mogging'' especially in the context of saying someone
                is ``mogging.''
            \item Draw a picture of your favorite animal.
        \end{enumerate}
    \end{questions}

\end{document}
